% Source: https://artofproblemsolving.com/community/c3h1950089_something_i_noticed
% Stable identity: aops-topic-1950089-q01
Write $n=\overline{ab}$ and $m=\operatorname{rev}(n)=\overline{ba}$, with $a,b\ge1$. Since $n^2$ ends in the last digit of $b^2\ne0$, $\operatorname{rev}(n^2)$ has as many digits as $n^2$. So $n^2$ and $m^2$ must have the same number of digits, $3$ or $4$.

\emph{Three digits.} Then $n\le31$ and $m\le31$, so $a,b\le3$, and $\overline{ab}\le31$, $\overline{ba}\le31$ leave $(a,b)\in\{(1,1),(1,2),(1,3),(2,1),(2,2),(3,1)\}$. In all these cases $a^2$, $2ab$, $b^2$ are at most $9$, so squaring involves no carries: $n^2$ has digits $a^2,2ab,b^2$ and $m^2$ has digits $b^2,2ab,a^2$, and the property holds. For example $13^2=169$, $31^2=961$; $22^2=484$.

\emph{Four digits.} Then $n,m\ge32$, so $a,b\ge3$. The leading digit of $\operatorname{rev}(n^2)$ is the last digit of $n^2$, which is the last digit of $b^2$, so the property forces: the leading digit of $n^2$ equals the last digit of $m^2$, which is the last digit of $a^2$. Check each $a$ from $3$ to $9$ (with $b\ge3$): the leading digit of $n^2$ for $n=\overline{a3},\dots,\overline{a9}$ is
\[
a=3{:}\ 1;\quad a=4{:}\ 1\text{ or }2;\quad a=5{:}\ 2,3;\quad a=6{:}\ 3,4;\quad a=7{:}\ 5,6;\quad a=8{:}\ 6,7;\quad a=9{:}\ 8,9,
\]
while the last digit of $a^2$ is $9,6,5,6,9,4,1$ respectively. They never agree, so there are no four-digit cases.

The numbers are $\boxed{11,\ 12,\ 13,\ 21,\ 22,\ 31}$.
